%
% Chapter 8 - Conclusion
%
\chapter{Conclusion} \label{cap:conclusion}

At the conclusion of the project timeline, EduCal has moved from a conceptual set of requirements to a functional web application that replaces the spreadsheet-based process previously used for institutional academic planning at ISEL (Chapter~\ref{cap:background}). Of the six functional requirements defined in Section~\ref{sec:functional-requirements}, four were fully implemented and are in use: hierarchical event and occurrence management (RF01, Sections~\ref{sec:core-entities} and~\ref{sec:frontend-event-dialogs}), the idempotent, rule-aware academic-year Migration Engine (RF02, Section~\ref{sec:migration-engine}), the configurable Business Rule Engine (RF03, Section~\ref{sec:rule-engine}), and role-based multi-user access control (RF06, Sections~\ref{sec:backend-authentication}, \ref{sec:backend-rbac}, and~\ref{sec:frontend-user-mgmt}). Two requirements from the original scope were not delivered: notification templates (RF04) and PDF schedule export (RF05); no corresponding functionality exists in the current codebase, and both are excluded from this delivery.

The decision to structure the backend as a decoupled Controller-Service-Data pipeline, with all storage access hidden behind a repository interface (Section~\ref{sec:data-access}), proved effective throughout the project. It allowed the migration engine and the rule engine to be built and iterated on independently of the underlying data store, and it made the transition from the JSON-backed mock data used while developing the business logic to the production PostgreSQL database a matter of changing a single environment variable, with no changes required to the controllers, services, API routes, or frontend. The same separation is what makes it possible to keep the lightweight Redis-backed implementation alongside PostgreSQL for development and demos, without either implementation leaking into the rest of the system.

Overall, EduCal delivers a reliable, working solution for the core problem stated in Chapter~\ref{cap:background} -- eliminating the manual, error-prone, spreadsheet-based rebuilding of the academic calendar every year -- through automated migration and configurable rule enforcement, backed by a secure, role-based multi-user system. Completing notification templates and schedule export, and moving from the current local/Vercel setup to the planned ISEL VM deployment (Chapter~\ref{cap:deployment}), remain as the identified gaps between this delivery and the full original scope.
